% ============================================================
%  J5 — APRÈS-MIDI : Crab Age Regression Challenge
%  Julien Rolland — M2 Développement Fullstack
% ============================================================
\documentclass[aspectratio=169, 10pt]{beamer}
\input{../shared/preamble}

\title[Challenge Kaggle]{Crab Age Regression Challenge}
\subtitle{Jour 5 --- Après-midi}
\date{Jour 5}

% ============================================================
\begin{document}
% ============================================================

\begin{frame}
  \titlepage
\end{frame}

\begin{frame}{Plan du module}
  \tableofcontents
\end{frame}

% ============================================================
\section{Le défi}
% ============================================================

\begin{frame}{La Mission : Expertise en Aquaculture}
  \begin{columns}[T]
    \begin{column}{0.52\textwidth}
      \begin{alertblock}{Le problème terrain}
        Déterminer l'âge des crabes est essentiel pour la gestion
        durable des stocks, mais l'analyse en laboratoire
        (section des carapaces) est \textbf{longue et coûteuse}.
        \begin{itemize}
          \item Processus invasif, impossible à grande échelle
          \item Expertise humaine rare et coûteuse
        \end{itemize}
      \end{alertblock}
      \smallskip
      \begin{block}{L'objectif}
        Prédire l'\textbf{âge (en mois)} d'un crabe à partir de
        mesures physiques non-invasives.
      \end{block}
    \end{column}
    \begin{column}{0.44\textwidth}
      \begin{exampleblock}{Votre rôle}
        En tant qu'ingénieurs IA, livrer le modèle le plus précis
        possible pour \textbf{automatiser cette tâche} sur le terrain.
        \vspace{4pt}
        \begin{center}
          \begin{tikzpicture}[font=\scriptsize,
            nd/.style={draw, rounded corners=3pt, align=center,
                       minimum width=2cm, minimum height=0.6cm,
                       fill=jedy_light!30}]
            \node[nd] (m) at (0,0)    {Mesures\\physiques};
            \node[nd] (ai) at (0,-1.1) {Modèle IA};
            \node[nd] (a) at (0,-2.2) {Âge estimé};
            \draw[->, thick] (m) -- (ai);
            \draw[->, thick] (ai) -- (a);
          \end{tikzpicture}
        \end{center}
      \end{exampleblock}
    \end{column}
  \end{columns}
\end{frame}

\begin{frame}{Le Dataset : Mesures de Crabes}
  \begin{columns}[T]
    \begin{column}{0.52\textwidth}
      \begin{block}{Features}
        \begin{itemize}
          \item \textbf{Sex :} Masculin (M), Féminin (F), Indéterminé (I)
          \item \textbf{Dimensions :} Longueur, Diamètre, Hauteur (mm)
          \item \textbf{Poids :} Total, sans coquille, viscères, carapace (g)
        \end{itemize}
      \end{block}
      \smallskip
      \begin{block}{Target}
        Colonne \textbf{Age} --- valeur continue (en mois).
      \end{block}
    \end{column}
    \begin{column}{0.44\textwidth}
      \begin{exampleblock}{Origine des données}
        Données synthétiques générées par IA.\\[4pt]
        Elles imitent les relations biologiques réelles tout en
        garantissant l'absence de solutions pré-existantes
        sur le web.
      \end{exampleblock}
      \smallskip
      \begin{alertblock}{Attention}
        Variable catégorielle \textbf{Sex} à encoder.\\
        Réfléchissez à la normalisation des poids.
      \end{alertblock}
    \end{column}
  \end{columns}
\end{frame}

% ============================================================
\section{Règles et logistique}
% ============================================================

\begin{frame}{La Métrique : Mean Absolute Error}
  \begin{columns}[T]
    \begin{column}{0.52\textwidth}
      \begin{block}{Définition}
        \[
          \mathrm{MAE} = \frac{1}{n}\sum_{i=1}^{n}|\hat{y}_i - y_i|
        \]
        C'est l'erreur moyenne en mois entre l'âge prédit et
        l'âge réel.
      \end{block}
      \smallskip
      \begin{exampleblock}{Traduction métier}
        «~Mon modèle se trompe en moyenne de $X$ mois
        sur l'âge du crabe.~»
      \end{exampleblock}
    \end{column}
    \begin{column}{0.44\textwidth}
      \begin{block}{MAE vs MSE}
        \begin{itemize}
          \item MSE pénalise \textbf{lourdement} les gros écarts
                (erreur quadratique)
          \item MAE est \textbf{robuste} aux outliers
          \item MAE est \textbf{directement interprétable}
                en unité métier
        \end{itemize}
      \end{block}
      \smallskip
      \begin{alertblock}{Objectif minimal}
        Battre la \textbf{baseline fournie} : arbre de décision
        simple --- MAE $\approx 1.45$ mois.
      \end{alertblock}
    \end{column}
  \end{columns}
\end{frame}

\begin{frame}{Logistique Kaggle}
  \begin{columns}[T]
    \begin{column}{0.52\textwidth}
      \begin{block}{Fichiers fournis}
        \begin{itemize}
          \item \textbf{train.csv :} données d'entraînement avec
                la colonne Age (15\,000 exemples)
          \item \textbf{test.csv :} données sans étiquettes ---
                vos prédictions ici (10\,000 exemples)
        \end{itemize}
      \end{block}
      \smallskip
      \begin{block}{Format de soumission}
        Fichier \textbf{submission.csv} avec deux colonnes :
        \texttt{id} et \texttt{Age}.\\[4pt]
        Sauvegardez dans \texttt{/kaggle/working/}.
      \end{block}
    \end{column}
    \begin{column}{0.44\textwidth}
      \begin{alertblock}{Leaderboard : Public vs Private}
        \begin{itemize}
          \item \textbf{Public (30\%) :} visible immédiatement
                --- votre score actuel
          \item \textbf{Private (70\%) :} révélé au J6 ---
                \textbf{c'est ce score qui compte}
          \item Attention à l'\textbf{overfitting} sur le
                leaderboard public !
        \end{itemize}
      \end{alertblock}
      \smallskip
      \begin{exampleblock}{Fin des soumissions}
        \textbf{Dimanche 22 mars à 23h55.}\\
        Votre validation locale est votre meilleur guide.
      \end{exampleblock}
    \end{column}
  \end{columns}
\end{frame}

% ============================================================
\section{Stratégie}
% ============================================================

\begin{frame}{Vos Leviers d'Action}
  \begin{columns}[T]
    \begin{column}{0.50\textwidth}
      \begin{block}{Prétraitement \& Feature Engineering}
        \begin{itemize}
          \item Encoder \textbf{Sex} : one-hot ou ordinal ?
          \item Normaliser les features continues ?
          \item Créer des \textbf{ratios} biologiques :\\
                ex. Poids / Longueur, Poids / Volume
          \item Gérer les valeurs aberrantes
        \end{itemize}
      \end{block}
      \smallskip
      \begin{block}{Validation locale}
        Créez votre propre split ou K-Fold.\\
        Ne touchez au test set qu'en soumission.
      \end{block}
    \end{column}
    \begin{column}{0.46\textwidth}
      \begin{block}{Choix du modèle}
        \begin{itemize}
          \item \textbf{Baseline :} arbre de décision (fourni)
          \item \textbf{Ensembles :} Random Forest,
                Gradient Boosting
          \item \textbf{Boosting :} XGBoost, LightGBM
          \item \textbf{MLP :} réseau de neurones simple
        \end{itemize}
      \end{block}
      \smallskip
      \begin{exampleblock}{Conseil}
        Le temps passé à comprendre vos données et créer
        des features pertinentes est souvent plus rentable
        que tester 50 algorithmes différents.
      \end{exampleblock}
    \end{column}
  \end{columns}
\end{frame}

% ============================================================
\section{Évaluation}
% ============================================================

\begin{frame}{Évaluation \& Notation (20\,\% du module)}
  \begin{columns}[T]
    \begin{column}{0.52\textwidth}
      \begin{block}{Performance Technique --- 10 points}
        \begin{itemize}
          \item Rang final sur le \textbf{Private Leaderboard}
          \item Condition de validation : battre la baseline
                «~prédire la moyenne~»
          \item Classement relatif entre groupes
        \end{itemize}
      \end{block}
      \smallskip
      \begin{block}{Soutenance Flash J6 --- 10 points}
        5 minutes par groupe :
        \begin{enumerate}
          \item Pipeline de préparation des données
          \item Modèles testés et choix final
          \item \textbf{Vos échecs} : qu'est-ce qui n'a pas marché ?
          \item Comment avez-vous lutté contre l'overfitting ?
        \end{enumerate}
      \end{block}
    \end{column}
    \begin{column}{0.44\textwidth}
      \begin{alertblock}{Ce qui différencie les équipes}
        \begin{itemize}
          \item Une vraie démarche expérimentale
          \item La rigueur de la validation locale
          \item La capacité à \textbf{expliquer} les choix
        \end{itemize}
      \end{alertblock}
      \smallskip
      \begin{exampleblock}{Critère clé de la soutenance}
        Les \textbf{échecs documentés} sont valorisés autant
        que les succès. L'ingénieur IA sait expliquer
        \textit{pourquoi} une piste n'a pas fonctionné.
      \end{exampleblock}
    \end{column}
  \end{columns}
\end{frame}

% ============================================================
\section*{Lancement}
% ============================================================

\begin{frame}{C'est parti !}
  \begin{columns}[T]
    \begin{column}{0.52\textwidth}
      \begin{block}{Accès à la compétition}
        \begin{enumerate}
          \item Connectez-vous sur \textbf{kaggle.com}
          \item Rejoignez via le lien d'invitation :\\
                {\footnotesize\texttt{kaggle.com/t/e57371844bcb4a0d96affb7875ec7373}}
          \item Créez un notebook Kaggle (GPU off, internet on)
          \item Exécutez la baseline --- première soumission
        \end{enumerate}
      \end{block}
      \smallskip
      \begin{block}{Organisation suggérée}
        \begin{itemize}
          \item \textbf{J5 PM :} EDA + baseline + feature engineering
          \item \textbf{J6 :} modèles avancés, tuning,
                soumission finale avant le \textbf{22/03 23h55}
        \end{itemize}
      \end{block}
    \end{column}
    \begin{column}{0.44\textwidth}
      \begin{exampleblock}{Par où commencer ?}
        \begin{enumerate}
          \item Ouvrir le \textbf{notebook baseline} fourni
          \item L'exécuter et soumettre --- score de référence
          \item Explorer les données : \texttt{.describe()}, corrélations
          \item Itérer : features, modèle, hyperparamètres
        \end{enumerate}
      \end{exampleblock}
      \smallskip
      \begin{alertblock}{Bonne chance !}
        Vous avez tous les outils vus en cours ---
        maintenant c'est à vous de jouer.
      \end{alertblock}
    \end{column}
  \end{columns}
\end{frame}

\end{document}
